% !TeX encoding = UTF-8
\documentclass[variant=apuntes,cover=institutional,mode=screen]{ua-engineering-v6-article}
% !TeX encoding = UTF-8
\UASetUniversity{Universidad Austral}
\UASetFaculty{Facultad de Ingeniería}
\UASetDepartment{Departamento de Computación}
\UASetCampus{Pilar}
\UASetCareer{Ingeniería Informática}
\UASetCommission{Comisión A}
\UASetProfessor{Equipo docente}
\UASetSemester{Segundo cuatrimestre 2026}
\UASetCourse{Sistemas Distribuidos}
\UASetDate{2026}


\UASetClassNumber{10}
\UASetClassTitle{Observabilidad distribuida}
\title{Clase 10 --- Logs, métricas, trazas y debugging causal en producción}
\author{\UAProfessor}
\date{\UADate}

\begin{document}
\maketitle

\section{Propósito de la clase}

La arquitectura distribuida multiplica fronteras, procesos, relojes y modos de falla. Una request puede atravesar servicios sincrónicos, topics, consumers y bases diferentes. Cuando algo falla, ya no existe un stack trace único ni una memoria compartida que cuente toda la historia.

Observabilidad es la capacidad de formular preguntas sobre el estado interno del sistema a partir de sus señales externas. No es sinónimo de recolectar datos. Un sistema puede emitir terabytes de telemetría y seguir siendo difícil de explicar si las señales carecen de contexto, correlación o relación con objetivos del usuario.

\begin{uakeybox}
La observabilidad útil convierte síntomas en hipótesis y las hipótesis en decisiones operativas verificables.
\end{uakeybox}

Los resultados esperados son:
\begin{itemize}
\item distinguir monitoreo de observabilidad;
\item diseñar logs estructurados con contexto seguro;
\item elegir métricas orientadas a servicio y recursos;
\item interpretar percentiles, cardinalidad y sampling;
\item construir trazas y propagar contexto;
\item correlacionar logs, métricas y trazas en un debugging causal;
\item diseñar SLI, alertas y telemetría para sistemas sincrónicos y asincrónicos.
\end{itemize}

\section{Caso conductor: checkout degradado}

A las 21:10, los errores 500 suben del 0.2\% al 8\%. La latencia p99 alcanza 12 segundos. La API tiene CPU normal, Payments muestra latencia irregular y la base tiene CPU 45\%, pero su pool de conexiones está completo.

Preguntas:
\begin{itemize}
\item ¿el problema comenzó en la API, Payments o la base?;
\item ¿afecta todas las regiones o una versión?;
\item ¿los retries son causa, consecuencia o ambas?;
\item ¿qué señales permiten distinguir hipótesis?;
\item ¿cuál es la mitigación inmediata y cuál el fix sistémico?
\end{itemize}

Este caso acompañará toda la clase.

\section{Monitoreo y observabilidad}

\subsection{Monitoreo}

Responde preguntas conocidas mediante métricas, umbrales y alertas:
\begin{itemize}
\item ¿se incumple el SLO?;
\item ¿creció la tasa de errores?;
\item ¿el consumer lag supera el objetivo?;
\item ¿un recurso está saturado?
\end{itemize}

\subsection{Observabilidad}

Permite investigar preguntas no predefinidas:
\begin{itemize}
\item ¿por qué solo fallan requests de una región y versión?;
\item ¿qué dependencia introduce la cola larga?;
\item ¿qué retry duplicó el efecto?;
\item ¿qué cambio amplió el blast radius?
\end{itemize}

No son opuestos. El monitoreo detecta desviaciones; la observabilidad ayuda a explicarlas.

\section{Tres señales complementarias}

\begin{itemize}
\item Métricas: comportamiento agregado, económico de consultar y alertar.
\item Logs: eventos discretos con detalle local y contexto.
\item Trazas: causalidad y latencia de una operación distribuida.
\end{itemize}

El valor aparece al correlacionarlas. Una métrica identifica el período; una traza selecciona un recorrido anómalo; los logs explican datos y decisiones internas.

\section{Logs estructurados}

\subsection{De frases a eventos consultables}

Un mensaje “Error procesando pedido” no permite agrupar por servicio, versión, región o causa. Un log estructurado puede incluir:
\begin{itemize}
\item timestamp y nivel;
\item servicio, versión, región e instancia;
\item trace ID, span ID, request/event ID;
\item order ID, payment ID u otro aggregate ID;
\item intento, deadline y política de retry;
\item latencia, resultado y error code;
\item mensaje humano para detalle.
\end{itemize}

Los nombres y tipos de campos deben ser estables. Si cada equipo usa nombres distintos, la correlación y consulta se degradan.

\subsection{Contexto y seguridad}

No todo dato debe loguearse. Tokens, contraseñas, números de tarjeta, secretos y PII innecesaria deben excluirse, enmascararse o tratarse según política. La observabilidad también es una superficie de seguridad y cumplimiento.

\subsection{Niveles}

\begin{itemize}
\item DEBUG: diagnóstico detallado, normalmente muestreado o temporal.
\item INFO: eventos relevantes de operación normal.
\item WARN: degradación o condición recuperable que merece atención.
\item ERROR: operación fallida que afecta una garantía o requiere investigación.
\end{itemize}

Si todo es ERROR, nada es accionable. Si los errores esperables se registran con stack trace masivo, el costo y ruido crecen.

\section{El pipeline de logs también falla}

Cloudflare informó un incidente en noviembre de 2024 donde, durante aproximadamente 3.5 horas, cerca del 55\% de los logs que normalmente enviaba a clientes no se entregaron y se perdieron.

Lecciones:
\begin{itemize}
\item ausencia de logs no demuestra ausencia de errores;
\item la completitud del pipeline de observabilidad necesita métricas y SLOs;
\item buffering, backpressure y retención deben diseñarse;
\item el observability stack forma parte del sistema crítico.
\end{itemize}

\section{Métricas}

\subsection{Qué comprimen}

Una métrica reduce muchos eventos a una serie temporal. Esa compresión la hace barata y útil, pero elimina detalle. Por eso debe responder una pregunta clara.

\subsection{Golden Signals}

Google SRE propone cuatro señales para servicios orientados a usuarios:
\begin{itemize}
\item latency;
\item traffic;
\item errors;
\item saturation.
\end{itemize}

Deben distinguirse solicitudes exitosas y fallidas: un error rápido no compensa una operación exitosa lenta.

\subsection{RED y USE}

RED para servicios:
\begin{itemize}
\item Rate;
\item Errors;
\item Duration.
\end{itemize}

USE para recursos:
\begin{itemize}
\item Utilization;
\item Saturation;
\item Errors.
\end{itemize}

RED muestra impacto de servicio; USE ayuda a investigar capacidad física. La CPU normal no descarta saturación de conexiones, locks, I/O o colas.

\section{Percentiles y cola larga}

Supongamos 95 requests de 100 ms y 5 requests de 8 s. La media puede parecer moderada, pero el 5\% de usuarios experimenta una demora extrema. Los percentiles describen la distribución:
\begin{itemize}
\item p50: experiencia típica;
\item p95: cola relevante;
\item p99/p99.9: extremos, a menudo dominados por fan-out y retries.
\end{itemize}

En una request que llama a muchos servicios, la probabilidad de tocar al menos una cola lenta aumenta. Por eso, el p99 de una dependencia importa aunque su media sea baja.

\section{Cardinalidad}

Una serie temporal suele identificarse por nombre y etiquetas. Si una etiqueta toma millones de valores, se crean millones de series.

Etiquetas razonables:
\begin{itemize}
\item endpoint normalizado;
\item status class;
\item region;
\item service version;
\item operation type.
\end{itemize}

Etiquetas peligrosas:
\begin{itemize}
\item request ID;
\item user ID;
\item raw URL con IDs;
\item error message libre;
\item stack trace.
\end{itemize}

Regla: identidades individuales viven mejor en logs y traces; las métricas agrupan dimensiones acotadas.

\section{Tracing distribuido}

\subsection{Trace y span}

Una trace representa el recorrido causal de una operación. Un span registra una unidad de trabajo con:
\begin{itemize}
\item inicio y duración;
\item parent span;
\item atributos;
\item eventos;
\item estado/resultado;
\item contexto de recursos.
\end{itemize}

La estructura no es solo cronológica. Parent/child explica causalidad y fan-out.

\subsection{Context propagation}

OpenTelemetry define propagación de contexto para conectar spans a través de procesos y redes. El contexto puede viajar en headers HTTP/gRPC o metadata de mensajes.

Si un servicio no extrae o inyecta el contexto, la traza se corta. Un consumidor asincrónico debe decidir si continúa la trace original, crea un vínculo o inicia una nueva trace correlacionada, según semántica y duración.

\subsection{Atributos}

Atributos útiles:
\begin{itemize}
\item operación y ruta normalizada;
\item peer service;
\item status code;
\item messaging destination;
\item retry attempt;
\item DB system y operation;
\item aggregate ID si la política lo permite.
\end{itemize}

Evitar payloads enormes y datos sensibles.

\section{Sampling}

Capturar todas las trazas puede ser costoso.

\subsection{Head sampling}
Decide al inicio. Es barato y distribuible, pero puede descartar incidentes raros antes de saber que fallarán.

\subsection{Tail sampling}
Decide después de observar la trace. Puede conservar errores, alta latencia o rutas específicas, pero necesita buffering y una infraestructura de decisión.

\subsection{Estrategia}

Una política razonable combina muestra base, retención total de errores/rutas críticas y adaptación durante incidentes. Debe evaluarse si permite responder preguntas reales.

\section{Correlación de señales}

Un workflow de investigación:
\begin{enumerate}
\item confirmar impacto mediante SLI/SLO;
\item delimitar período, región, versión, endpoint o tenant;
\item revisar cambios recientes;
\item observar RED/Golden Signals;
\item comparar recursos con USE;
\item seleccionar traces lentas/fallidas y normales;
\item localizar el primer span degradado;
\item consultar logs con trace ID;
\item formular hipótesis y buscar evidencia que podría refutarla;
\item mitigar, medir y documentar.
\end{enumerate}

Este método evita saltar directamente a una explicación favorita.

\section{Síntoma y causa sistémica}

En el caso conductor:
\begin{itemize}
\item síntoma: API 500 y p99 12s;
\item causa próxima: pool de conexiones agotado;
\item cadena: query lenta o deadline largo $\rightarrow$ requests retenidas $\rightarrow$ timeouts $\rightarrow$ retries $\rightarrow$ más conexiones $\rightarrow$ saturación.
\end{itemize}

Mitigar puede significar reducir retries, abrir circuit breaker, aumentar capacidad temporal o revertir configuración. El fix sistémico puede requerir query, timeout budget, pool isolation y load shedding.

\section{Caso real: propagación global de una configuración}

Cloudflare publicó un postmortem del 18 de noviembre de 2025. Un cambio de permisos hizo que una consulta produjera múltiples entradas en un archivo usado por Bot Management; el archivo duplicó su tamaño y fue propagado a máquinas de la red, provocando un outage amplio.

Lecciones para observabilidad:
\begin{itemize}
\item tamaño y forma de un artefacto son métricas operativas;
\item rollout global necesita canaries y límites;
\item el disparador local no explica el blast radius por sí solo;
\item postmortems conectan cambio, propagación, defensas y detección.
\end{itemize}

\section{Sistemas asincrónicos y sagas}

La causalidad no siempre es una cadena. Un evento puede activar varios consumidores y una saga puede durar minutos.

Señales útiles:
\begin{itemize}
\item event ID, correlation ID y saga ID;
\item lag y edad del mensaje;
\item estado de saga y paso actual;
\item retries, DLQ y redrive;
\item tiempo desde evento inicial hasta efecto final;
\item cantidad de compensaciones.
\end{itemize}

Una trace tradicional puede usar links entre producer y consumers en vez de un parent único.

\section{Alertas y SLOs}

Las alertas útiles se orientan a síntomas de usuario o consumo rápido del error budget. Ejemplos:
\begin{itemize}
\item checkout success rate bajo objetivo;
\item p99 por encima del SLO;
\item pipeline con edad de evento mayor a límite;
\item disponibilidad de una capacidad crítica.
\end{itemize}

Alertas de CPU o disco pueden apoyar diagnóstico, pero no deberían pagear sin una relación accionable con el servicio, salvo riesgo inmediato de agotamiento.

\section{Telemetría como producto interno}

Para cada señal:
\begin{itemize}
\item definir pregunta y usuario;
\item schema y ownership;
\item volumen, cardinalidad y costo;
\item retención;
\item privacidad;
\item completitud y calidad;
\item dependencia de operación;
\item política de degradación si el pipeline falla.
\end{itemize}

La observabilidad tiene presupuesto. Debe priorizar caminos críticos y preguntas valiosas.

\section{Método de diseño}

\begin{enumerate}
\item definir journeys críticos y SLOs;
\item instrumentar límites de servicio y dependencias;
\item estandarizar contexto y campos;
\item diseñar métricas con cardinalidad acotada;
\item propagar traces en sync y async;
\item definir sampling y retención;
\item diseñar dashboards por preguntas;
\item alertar por impacto y burn rate;
\item ensayar fallas para comprobar evidencia;
\item revisar costo y utilidad periódicamente.
\end{enumerate}

\section{Errores frecuentes}

\begin{itemize}
\item confundir telemetría con observabilidad;
\item logs de texto sin contexto;
\item datos sensibles en logs;
\item métricas con cardinalidad explosiva;
\item promedios sin percentiles;
\item traces cortadas;
\item sampling que descarta errores raros;
\item dashboards sin pregunta;
\item alertas por todo;
\item investigar sin hipótesis refutable;
\item asumir que ausencia de señal equivale a salud.
\end{itemize}

\section{Bibliografía guiada y casos reales}

\begin{itemize}
\item Google SRE Book, \href{https://sre.google/sre-book/monitoring-distributed-systems/}{Monitoring Distributed Systems}: Golden Signals y principios de monitoreo.
\item OpenTelemetry, \href{https://opentelemetry.io/docs/concepts/signals/traces/}{Traces} y \href{https://opentelemetry.io/docs/concepts/context-propagation/}{Context Propagation}.
\item Kleppmann, capítulos sobre fallas, logs y sistemas derivados.
\item Cloudflare, \href{https://blog.cloudflare.com/cloudflare-incident-on-november-14-2024-resulting-in-lost-logs/}{incidente de pérdida de logs de noviembre de 2024}.
\item Cloudflare, \href{https://blog.cloudflare.com/18-november-2025-outage/}{postmortem del outage del 18 de noviembre de 2025}.
\end{itemize}

\section{Conexión con el Trabajo Final Integrador}

Cada grupo debe definir:
\begin{itemize}
\item tres journeys críticos;
\item SLI/SLO;
\item logs estructurados y campos prohibidos;
\item métricas RED/USE;
\item propagación de trace en RPC y eventos;
\item dashboard y alertas;
\item una investigación de falla simulada;
\item costo/retención de telemetría.
\end{itemize}

\section{Autoevaluación}

\begin{enumerate}
\item ¿Qué diferencia monitoreo de observabilidad?
\item ¿Qué campo conecta logs de una misma request?
\item ¿Por qué un error message es mala etiqueta de métrica?
\item ¿Qué muestran p95/p99 que la media oculta?
\item ¿Qué diferencia RED y USE?
\item ¿Qué es un span huérfano?
\item ¿Cuándo tail sampling es valioso?
\item ¿Cómo se observa una saga asincrónica?
\item ¿Por qué el pipeline de logs necesita SLO?
\item ¿Qué evidencia refutaría su hipótesis de causa raíz?
\end{enumerate}

\section{Síntesis final}

\begin{uakeybox}
Las métricas encuentran el patrón, las trazas reconstruyen el recorrido y los logs explican el detalle. La observabilidad aparece cuando esas señales comparten contexto, se orientan a objetivos y permiten formular y refutar hipótesis sobre el sistema real.
\end{uakeybox}

\end{document}
